\section{Poisson differences and reference coefficients}
\label{p2:poisson-calculus}

\subsection{Computing the correction coefficient}
\label{p1:discussion:charlier}

The coefficient in the total variation equivalent can be
evaluated from four Poisson masses.  We record the elementary
Poisson--Charlier calculation, together with a bound for every
order. These are standard identities for signed Poisson
approximations. The orthogonality bound and the exact
second-difference norm appear in Zacharovas and
Hwang~\cite[(2.13), Remark~4.4]{zacharovashwang2010}; see also
\cite{bkn2014} for the Poisson--Charlier expansion framework.
We include the calculation to fix the normalization of
\(\Delta\), the endpoint conventions, and the factor relating
the \(\ell^1\) norm to total variation.

\begin{proposition}[Norms of Poisson differences]
\label{p1:discussion:charlier-norm}
For $\mu>0$, put $\pi_\mu(k)=e^{-\mu}\mu^k/k!$ when $k\ge0$
and $\pi_\mu(k)=0$ when $k<0$.  Define
\[
 u_\pm=\mu+\frac12\pm\frac{\sqrt{4\mu+1}}2,\qquad
 a=\lfloor u_-\rfloor+1,\qquad b=\lceil u_+\rceil-1.
\]
Then
\begin{equation}
 \|\Pi_\mu*\Delta^{*2}\|_1
 =2\bigl\{\pi_\mu(a-1)-\pi_\mu(a-2)
             +\pi_\mu(b-1)-\pi_\mu(b)\bigr\}.
 \label{p1:discussion:finite-norm}
\end{equation}
More generally, for every integer $r\ge0$,
\begin{equation}
 \|\Pi_\mu*\Delta^{*r}\|_1
       \le\min\left\{2^r,\frac{\sqrt{r!}}{\mu^{r/2}}\right\}.
 \label{p1:discussion:general-norm}
\end{equation}
In particular,
\begin{equation}
 \|\Pi_2*\Delta^{*2}\|_1=\frac{10}{3e^2},
 \qquad \mathcal C(2)=\frac{5}{6e^2}.
 \label{p1:discussion:lambda-two}
\end{equation}
\end{proposition}

The proof, including the normalization of the Charlier polynomials and
the telescoping at integer roots, is given in
Section~\ref{p14:comp:poisson-differences}.

The function $\mathcal C$ is continuous and strictly positive
on $(0,\infty)$.  Indeed, coupling two Poisson variables by
an independent increment gives
$\|\Pi_\mu-\Pi_\nu\|_1\le2(1-e^{-|\mu-\nu|})$.
Convolving with $\Delta^{*2}$ proves continuity; the mass
at zero gives positivity, as in
\eqref{p1:close:constant-bounds}.  Thus the finite formula
remains continuous when one of its two roots crosses an integer.

For completeness, let
$H_{\mu,\theta}=\Law(\Pois(\mu-\theta)+2\Pois(\theta/2))$,
with independent summands and $0\le\theta<\mu$.
The expansion of its convolution exponential gives the
parameter-uniform estimate
\begin{equation}
 \left|\TV(H_{\mu,\theta},\Pi_\mu)
                     -\theta\mathcal C(\mu)\right|
       \le\theta^2e^{2\theta}.
 \label{p1:discussion:compound-coefficient}
\end{equation}
To verify the constant, write
$H_{\mu,\theta}
 =\Pi_\mu*\exp_*(\theta\Delta^{*2}/2)$.
The $\ell^1$ norm of the exponent is at most $2\theta$,
so the remainder after its linear term has norm at most
$2\theta^2e^{2\theta}$.
The reverse triangle inequality and
$\TV=\|\cdot\|_1/2$ prove
\eqref{p1:discussion:compound-coefficient}.
In particular, the factor $1/4$ in $\mathcal C$ consists
of the factor $1/2$ in the quadratic correction and the
factor $1/2$ in the total variation convention.
Neither \eqref{p1:discussion:general-norm} nor the remainder
constant in \eqref{p1:discussion:compound-coefficient} is
asserted to be optimal.


\subsection{Norms of Poisson differences}
\label{p14:comp:poisson-differences}

The following calculation proves Proposition~\ref{p1:discussion:charlier-norm}.
It fixes both the signed-difference convention and the factor of two between
$\ell^1$ distance and total variation.

\begin{proof}[Proof of Proposition~\ref{p1:discussion:charlier-norm}]
Write $v_k=(\Pi_\mu*\Delta^{*2})\{k\}$.  Direct calculation gives
\[
 v_k=\pi_\mu(k)-2\pi_\mu(k-1)+\pi_\mu(k-2)
     =\frac{\pi_\mu(k)}{\mu^2}\bigl((k-\mu)^2-k\bigr).
\]
The quadratic has roots $u_-,u_+$ and is strictly negative
exactly at the integers $a,\ldots,b$.  These form a nonempty
interval; integer roots contribute zero and are excluded by
the definitions of $a,b$.  Since $\sum_kv_k=0$,
\[
 \begin{aligned}
 \sum_k|v_k|
   &=-2\sum_{k=a}^b v_k\\
   &=2\bigl\{\pi_\mu(a-1)-\pi_\mu(a-2)
                 +\pi_\mu(b-1)-\pi_\mu(b)\bigr\},
 \end{aligned}
\]
by telescoping.  This proves \eqref{p1:discussion:finite-norm}.
For $\mu=2$, the quadratic is $(k-1)(k-4)$: the negative
masses are $v_2=-e^{-2}$ and $v_3=-2e^{-2}/3$.
Their sum, multiplied by $-2$, gives
\eqref{p1:discussion:lambda-two}.  The same calculation at
$\mu=1$ gives the value $3/e$ used in
\eqref{p1:close:dyadic-example}.

For the general bound, define the normalized Charlier polynomials
\[
 C_r(k;\mu)=\sum_{j=0}^r
       \binom rj(-1)^{r-j}\frac{(k)_j}{\mu^j}.
\]
Then $(\Pi_\mu*\Delta^{*r})\{k\}=\pi_\mu(k)C_r(k;\mu)$.
Their exponential generating function is
\[
 \sum_{r\ge0}C_r(k;\mu)\frac{t^r}{r!}
       =e^{-t}(1+t/\mu)^k.
\]
For $P\sim\Pois(\mu)$, the expectation of the product of
the generating functions with variables $s,t$ is
\[
 \E\!\left[e^{-s-t}(1+s/\mu)^P(1+t/\mu)^P\right]
       =e^{st/\mu}.
\]
Coefficient comparison, justified in a neighborhood of
$(s,t)=(0,0)$ by the Poisson exponential moments, yields
\[
 \E[C_r(P;\mu)C_\ell(P;\mu)]
       =\ind_{\{r=\ell\}}\,\frac{r!}{\mu^r}.
\]
Cauchy--Schwarz now gives the second bound in
\eqref{p1:discussion:general-norm}.  The first follows from
$\|\Pi_\mu\|_1=1$, $\|\Delta\|_1=2$, and submultiplicativity
of the convolution norm.
\end{proof}
